\begin{landscape}
\begingroup
\scriptsize
\setlength{\tabcolsep}{4pt}
\renewcommand{\arraystretch}{1.16}

\section*{Supplementary Information}

\noindent\textbf{Supplementary Table S1.} Cohorts, sample counts, and feature
families used in the dCST applications.

\begin{longtable}{p{0.16\linewidth}p{0.38\linewidth}p{0.38\linewidth}}
\toprule
Analysis layer & Cohorts and sample counts & Feature family and filtering \\
\midrule
\endfirsthead
\toprule
Analysis layer & Cohorts and sample counts & Feature family and filtering \\
\midrule
\endhead
\bottomrule
\endfoot
Gut AGP discovery and exploratory prioritization &
Local AGP stool table derived from PRJEB11419; 24,566 source samples and 24,421 quality-controlled samples after library-size and all-zero post-filter exclusions. AGP phenotype fields were used for exploratory prioritization, not diagnosis-grade CD/UC inference. &
Shared-taxonomy stool taxon matrix. ASV sequences were reassigned with the DADA2/SILVA workflow used for the external cohorts. Retained taxa required count \(\geq 2\) in at least 5\% of retained samples, yielding 340 taxa. \\
\midrule
Gut external IBD transfer and rebuilt-cohort analysis &
Main external IBD cohorts: Halfvarson 2017 stool branch, 533 samples from 118 subjects for pooled IBD versus healthy; HMP2/IBDMDB mucosal slice, 166 samples from 81 subjects; Gevers 2014 Crohn-focused stool branch, 396 samples from 363 subjects. Halfvarson subtype branches contained 200 Crohn-case samples from 49 subjects and 279 UC-case samples from 60 subjects, each compared with 54 healthy-control samples from 9 subjects. &
External cohorts were aligned to the AGP retained-feature dictionary for direct transfer. Rebuilt analyses used cohort-specific filtered matrices under the same absorb workflow, allowing related dominance structure to be rediscovered without requiring exact AGP labels. \\
\midrule
CT full-VOG hierarchy, capped robustness runs, and functional follow-up &
The 3,602-sample CT+reference learning cohort contained 604 outcome-labeled CT metagenomes, 434 additional unlabeled CT metagenomes, and 2,564 VIRGO2 reference metagenomes. The 604 labeled CT samples were used for all CT outcome contrasts: 298 spontaneous-clearance and 306 persistence samples. &
Primary high-resolution feature family was the keep-filtered VIRGO2 VOG matrix. From 960,439 VOGs in the combined union table, VOGs were retained if observed in at least 10 samples and at least 1\% of the combined cohort, yielding 509,501 VOGs. Capped robustness runs used 25k, 50k, 100k, and 250k VOG subsets drawn from the same retained universe by variance ranking. \\
\midrule
Cross-application synthesis and assets &
The synthesis integrates the gut transfer/rebuilt analyses and CT full-VOG analyses as complementary evaluations of the dCST state vocabulary. &
No new feature family is introduced for the synthesis; it reuses the gut shared-taxonomy and CT full-VOG feature families above. \\
\end{longtable}

\clearpage

\noindent\textbf{Supplementary Table S2.} dCST construction, testing families,
and transfer rules used in the dCST applications.

\begin{longtable}{p{0.16\linewidth}p{0.34\linewidth}p{0.42\linewidth}}
\toprule
Analysis layer & dCST support schedule & Testing, BH family, and transfer rule \\
\midrule
\endfirsthead
\toprule
Analysis layer & dCST support schedule & Testing, BH family, and transfer rule \\
\midrule
\endhead
\bottomrule
\endfoot
Gut AGP discovery and exploratory prioritization &
Absorb-policy hierarchy fit through depth 4. Minimum support: \(n_0=50\) at depth 1; \(n_0=25\) at depths 2--4. Pure-policy results were retained as sensitivity analyses. &
AGP condition screens and label follow-up used depth-specific panel and label families with Benjamini--Hochberg correction. Bayesian follow-up was supportive; BH-adjusted frequentist results were treated as screen-level evidence. No external transfer is involved in this discovery row. \\
\midrule
Gut external IBD transfer and rebuilt-cohort analysis &
Direct transfer used the frozen AGP absorb hierarchy through depths 1--4. Rebuilt-cohort analyses used the same absorb-policy construction and depth range used for the AGP-derived hierarchy. &
Direct transfer used \texttt{transfer.dcsts()} with the corrected nearest-lineage-set rule: external samples were compared against realized AGP lineage sets at each nominal depth, and stable terminal lineages remained assignable at later nominal depths. Rebuilt mode analyzed each external cohort independently. Label-level and panel-level tests used BH correction within the prespecified comparison/mode/depth families. Repeated-subject sensitivity retained one sample per subject and used bootstrap one-sample-per-subject recurrence for leading rows. \\
\midrule
CT full-VOG hierarchy, capped robustness runs, and functional follow-up &
Absorb-policy full-VOG and capped-VOG hierarchies fit through depth 4. Minimum support: \(n_0=50\) at depth 1; \(n_0=25\) at depth 2; \(n_0=12\) at depth 3; \(n_0=10\) at depth 4. &
Primary readout was panel-level CT outcome association across the full dCST label panel at each depth, with BH correction over the prespecified depth-level family. Label-level Fisher tests and Jeffreys-prior Bayesian prevalence contrasts were supportive sparse-label summaries. Product, taxonomy, and KEGG contrasts were descriptive post hoc interpretation layers, not independent discovery gates. No independent CT replication cohort is included. \\
\midrule
Cross-application synthesis and assets &
No additional dCST hierarchy is fit for the synthesis figure. &
The synthesis is descriptive and does not introduce a new testing family. It summarizes the complementary application evidence: gut portability and external transfer, and CT biological and functional interpretability. \\
\end{longtable}

\clearpage

\noindent\textbf{Supplementary Table S3.} Implementation-level settings for
deterministic assignment, stochastic inferential summaries, and multiple-testing
families.

\begin{longtable}{p{0.24\linewidth}p{0.68\linewidth}}
\toprule
Implementation item & Setting used in the reported analyses \\
\midrule
\endfirsthead
\toprule
Implementation item & Setting used in the reported analyses \\
\midrule
\endhead
\bottomrule
\endfoot
dCST package-operation mapping &
The \texttt{linf} package functions used for Algorithm~\ref{alg:dcst-fit-transfer} were: \texttt{normalize.linf()} for sample-wise normalization; \texttt{linf.cells()} for dominant-feature assignment; \texttt{linf.csts()} for depth-1 dCST construction with pure or absorb low-support policies; \texttt{refine.linf.csts()} and \texttt{refine.linf.csts.iter()} for recursive depth-wise refinement, including absorb reassignment; and \texttt{transfer.dcsts()} for assigning aligned samples into a frozen fitted hierarchy. Application-specific repository workflows then performed cohort filtering, rebuilt-cohort analysis, association testing, and biological follow-up. \\
\midrule
dCST tie handling and feature identity &
All primary dCST pipelines used the \texttt{linf} default \texttt{tie.method = "first"}. For primary dominance assignment, tied features are resolved by the fixed feature-dictionary/column order; duplicate feature IDs or labels are made unique before assignment. Absorb reassignment uses the same ordered retained-feature dictionary restricted to the retained sibling targets for the active parent. If a low-support or all-zero sample has no positive retained-sibling value, it is assigned to the fallback retained sibling defined by that same ordered retained-target list. Observed top-feature ties were absent or rare in the reported audits, so tie handling is specified for reproducibility rather than because ties drive the analyses. \\
\midrule
Monte Carlo Fisher settings &
AGP exploratory condition-depth screening used the recorded seed \(20260411\), standard Monte Carlo Fisher mode, and \(B=1{,}000\) simulations per panel in the run reported here. CT full-VOG omnibus panels used Fisher tests with Monte Carlo evaluation and \(B=50{,}000\) simulations per panel. The CT association scripts did not set a separate Fisher-test seed, so exact simulated \(p\)-value digits are reproduced from the generated result tables rather than regenerated deterministically from the script alone. dCST labels themselves do not depend on Monte Carlo draws. \\
\midrule
Bayesian sparse-label summaries &
Gut follow-up used 4,000 posterior draws for Jeffreys-prior beta-binomial prevalence contrasts, \(p_{\mathrm{in}}\sim\mathrm{Beta}(a+0.5,b+0.5)\) and \(p_{\mathrm{out}}\sim\mathrm{Beta}(c+0.5,d+0.5)\), and used \texttt{arm::bayesglm} binomial logistic models for adjusted label-presence summaries when complete-case support was sufficient. Available AGP covariates were host age bin, sex, and BMI. CT follow-up used the same Jeffreys-prior prevalence contrast with 4,000 draws and did not run gut-style adjusted Bayesian logistic regression. CT Bayesian follow-up runs used seeds \(20260411\) and \(20260415\). \\
\midrule
Benjamini--Hochberg families &
AGP exploratory phenotype screening adjusted the eligible phenotype-by-depth omnibus panels as one screen family; AGP label follow-up reported both global label-level BH values across tested labels/depths within each outcome and depth-specific BH values. Gut external transfer and rebuilt-cohort analyses adjusted label tests within each prespecified cohort/comparison/mode/depth family. CT full-VOG and capped-VOG omnibus tests adjusted across depths within each feature-cap run; label-level tests adjusted within each depth panel. \\
\end{longtable}

\clearpage

\noindent\textbf{Supplementary Table S4.} Setting-dependence and robustness
checks for deterministic dCST construction.

\begin{longtable}{p{0.18\linewidth}p{0.48\linewidth}p{0.26\linewidth}}
\toprule
Check & Summary & Interpretation \\
\midrule
\endfirsthead
\toprule
Check & Summary & Interpretation \\
\midrule
\endhead
\bottomrule
\endfoot
Deterministic tie rule & All reported dCST fits used \texttt{tie.method = "first"}: first retained feature for primary dominance and first retained target for absorb reassignment. This is part of the label definition, not a random analysis option. & The reported labels are reproducible by definition once feature order is fixed. \\
\midrule
Observed top-feature ties & Dominant-feature ties were absent in the bundled AGP gut example (0/766 rows) and rare in the CT cap audit: 1/3602 rows in the 50k-VOG cap audit (0.03\%); maximum tied top set size 2. & Tie handling is specified for exact reproducibility; available audits do not suggest that top-feature ties drive the reported examples. \\
\midrule
Gut support-schedule perturbation & Using the existing AGP raw-lineage size audit, a \(\pm25\%\) perturbation around the manuscript support schedule 50/25/25/25 changed pre-policy hierarchy granularity but preserved high pre-absorb sample coverage at shallow depths. These diagnostic counts are pre-absorb raw-lineage counts and should not be compared directly with the post-absorb retained-state counts in Table~1. -25\%: thresholds 38/19/19/19; retained raw lineages 45/170/221/123; pre-absorb retained sample share 96.8\%/83.3\%/52.1\%/16.8\%. nominal: thresholds 50/25/25/25; retained raw lineages 40/141/154/89; pre-absorb retained sample share 96.2\%/81.3\%/47.4\%/14.4\%. +25\%: thresholds 63/31/31/31; retained raw lineages 35/115/114/64; pre-absorb retained sample share 95.3\%/79.0\%/43.8\%/12.1\%. & The support schedule controls naming granularity and should be treated as an explicit modeling choice, not as a hidden clustering parameter. \\
\midrule
CT retained-feature cap robustness & Capped VOG analyses used nested variance-ranked subsets from the same 509,501-VOG retained universe. 25k: depth-3 q=0.0414, V=0.325, depth-4 q=0.2125; 50k: depth-3 q=0.0053, V=0.360, depth-4 q=0.0074; 100k: depth-3 q=0.0042, V=0.358, depth-4 q=0.0063; 250k: depth-3 q=0.0044, V=0.357, depth-4 q=0.0053; 509,501: depth-3 q=0.0046, V=0.357, depth-4 q=0.0047. The 25k cap was shallower, whereas 50k and larger caps retained depth-4 support. & This supports the CT result as a within-cohort feature-family robustness result, not as external replication. \\
\midrule
Gut feature-filter scope & The gut taxonomic analysis is conditional on the stated shared-taxonomy retained family: count \(\geq2\) in at least 5\% of retained AGP samples, yielding 340 taxa. The combined manuscript does not claim robustness to all alternative gut prevalence filters. & This bounds the gut portability claim to the retained feature family used for transfer and rebuilt analyses. \\
\end{longtable}


\endgroup
\end{landscape}
